\documentclass[reqno]{amsart}
\usepackage[utf8]{inputenc}
\usepackage{amsmath,amssymb,amsthm}
\usepackage{hyperref}

\newtheorem{theorem}{Theorem}[section]
\newtheorem{lemma}[theorem]{Lemma}
\newtheorem{corollary}[theorem]{Corollary}
\newtheorem{proposition}[theorem]{Proposition}
\newtheorem{definition}[theorem]{Definition}
\newtheorem{remark}[theorem]{Remark}

\newcommand{\R}{\mathbb{R}}
\newcommand{\C}{\mathbb{C}}
\newcommand{\Q}{\mathbb{Q}}

\title[What is the classical Regge limit for $n\geq5$?]{What is the classical Regge limit for
$n\geq5$ particles? An open problem, with computational evidence from a five-photon degree-11
search}

\author{Mart\'in Baca}
\date{}

\begin{document}

\begin{center}
\fbox{\parbox{0.85\textwidth}{\centering
\textbf{Preprint v1.0 --- not peer-reviewed.} This note reports a negative/exploratory result and
an open problem, not a resolved theorem. Independently checked in-house (see
Section~\ref{sec:verification}), but not yet reviewed by anyone outside this project. Feedback,
corrections, and criticism are genuinely welcome --- please open an issue at
\url{https://github.com/martinbacaia/crg-open-problem-n5} or contact the author.}}
\end{center}

\begin{abstract}
The Classical Regge Growth (CRG) conjecture --- that any consistent classical $S$-matrix grows no
faster than $s^2$ in the Regge limit $s\to\infty$ at fixed $t$ --- is well-posed and has been used
to constrain four-particle photon and graviton $S$-matrices~\cite{cgghjm2019}. Extending this to
$n\geq5$ particles requires a generalization of ``the Regge limit'' itself, and we show this
generalization is genuinely ambiguous: no definition exists in the literature, and the closest
established framework with a similar name --- multi-Regge kinematics, from the gauge-theory
literature on multi-particle production~\cite{delduca1993} --- is motivated by different physics
and is not obviously the correct generalization of the causality-derived CRG bound. Using the
degree-9 and degree-11 five-photon operator basis constructed in a companion
note~\cite{fivephotons}, we construct three qualitatively distinct, explicitly validated families
of $n=5$ Regge trajectories and search for operators satisfying the conjectured bound. We find:
the unique degree-9 invariant has an exact growth floor of $s^3$ within one family; a
9-dimensional subspace of the (at-the-time) 15-dimensional degree-11 space satisfies the bound in
that same family; a 2-dimensional subspace survives two qualitatively different families and two
independent parameter configurations each; but a third family kills this subspace entirely, and an
exact computation of the combined kernel across all three families, on the complete 16-dimensional
degree-11 space (including a topology discovered after the two-family result), gives exactly
$\{0\}$. We interpret this pattern --- four structurally different attempts to enlarge the search
all failing in the same way --- as evidence that the obstruction is in the definition of the
limit, not in an incomplete search, and we state this as the open problem.
\end{abstract}

\maketitle
\tableofcontents

\section{Introduction}

The Classical Regge Growth (CRG) conjecture states that the tree-level $S$-matrix of any
consistent classical theory grows no faster than $s^2$ in the Regge limit ($s\to\infty$ at fixed
$t$)~\cite{cgghjm2019}. There is evidence for this from causality: the chaos bound in AdS/CFT,
applied to the boundary dual of a bulk theory of scalars, gauge fields, or gravitons, implies CRG
for the corresponding flat-space $S$-matrices~\cite{chaosbound2021}. Deriving CRG this way is
substantial holographic technical work, well beyond the scope of this project; instead, following
Chowdhury, Gadde, Gopalka, Halder, Janagal \& Minwalla (CGGHJM)~\cite{cgghjm2019}, we take CRG as a
given conjecture and ask what it implies for the operator classification at $n=5$, building on the
representation-theoretic and combinatorial foundation constructed in a companion
note~\cite{fivephotons}.

At $n=4$, ``the Regge limit'' has an unambiguous meaning: elastic $2\to2$ scattering at fixed
momentum transfer $t$, realized physically by boosting all four momenta while correlating the
scattering angle to keep $t$ fixed as $s\to\infty$. At $n\geq5$, there is no single kinematic
invariant playing the role of ``$t$'', and it is not obvious a priori which invariants should be
held fixed while which others diverge. This note's central finding is that this ambiguity is not
merely bookkeeping: it appears to be the actual obstruction to extending the $n=4$ classification
program to $n=5$, at least within the family of constructions we tried.

\begin{remark}[Literature search]
A direct search of the arXiv API (both for ``classical Regge growth'' specifically, and for
follow-up work on the $n=4$ classification) turned up no paper addressing CRG at $n\geq5$; the
follow-up literature since 2019 extends the $n=4$ story (massive spinning particles, dimensional
reduction, the holographic chaos-bound derivation itself) but does not touch $n\geq5$. Separately,
there is an established, decades-old, and extensively used framework called \emph{multi-Regge
kinematics} (MRK) for $n$-point amplitudes at $n\geq5$, going back to Del
Duca~\cite{delduca1993} and used throughout the BFKL/planar-$\mathcal{N}=4$ literature up to
$n=7$-point amplitudes: particles are ordered by rapidity, and the limit takes all consecutive
rapidity gaps to infinity (``strong rapidity ordering'') with transverse momenta fixed. We do
\emph{not} adopt MRK here, for a specific reason: MRK's physical motivation is multi-particle
production at high energy in gauge theory (an inherently different regime from the causality/chaos
bound that motivates CRG), and nothing we are aware of establishes that MRK is the kinematic
regime in which a chaos-bound-type causality argument would actually constrain growth to $s^2$.
Adopting it without that justification would repeat the same error as inventing an ad hoc
generalization, only with a citation that lends false confidence. We regard this as itself part of
the open problem, not a solved input.
\end{remark}

\subsection*{Organization} Section~\ref{sec:n4} reviews the $n=4$ Regge limit at the level of
explicit momenta, validated as a template. Section~\ref{sec:naive} shows a naive generalization to
$n=5$ fails structurally. Section~\ref{sec:families} constructs three qualitatively distinct
families of $n=5$ Regge trajectories. Section~\ref{sec:results} reports the growth rates found at
degrees 9 and 11, culminating in the negative result on the full 16-dimensional degree-11 space.
Section~\ref{sec:verification} summarizes the verification methodology, including several bugs
found and fixed. Section~\ref{sec:discussion} discusses what this means and states the open
problem.

\section{The \texorpdfstring{$n=4$}{n=4} Regge limit, validated at the level of explicit momenta}
\label{sec:n4}

We work in $D=6$ throughout (consistent with~\cite{fivephotons}), with the standard realization of
elastic $2\to2$ scattering at fixed $t$:
\[
p_1=-E(1,n_1),\quad p_2=-E(1,-n_1),\quad p_3=E(1,n_3),\quad p_4=E(1,-n_3),
\]
with $n_1,n_3$ unit spatial vectors at angle $\theta$, giving $s=4E^2\to\infty$ and
$t=-2E^2(1-\cos\theta)$. Choosing $\cos\theta(E)=1+t_0/(2E^2)$ keeps $t=t_0$ exactly fixed as
$E\to\infty$ ($\theta\to0$, correlated with $E\to\infty$) --- the standard forward-limit
realization of ``$s\to\infty$ at fixed $t$''. Crucially, all four momenta participate in the
high-energy limit; no particle is held at rest.

\begin{remark}[A gauge-reference bug, instructive for $n=5$]
\label{rem:gaugebug}
Constructing the polarizations $\varepsilon_i$ requires a reference momentum to fix the residual
gauge freedom (projecting a seed vector to satisfy $\varepsilon_i\cdot p_i=0$). An initial attempt
used a ``$t$-channel partner'' as that reference (e.g.\ $p_3$ for $\varepsilon_1$) and found
$(F^1{:}F^2)(F^3{:}F^4)$ growing as $\Lambda^5$, not the $\Lambda^4$ expected by naive degree
counting. Diagnosis: $\eta(p_1,p_3)=-t/2$ is held \emph{fixed} by construction (it is literally the
momentum transfer being kept constant), so using $p_3$ as a gauge reference divides by a quantity
that does not grow, causing the transverse part of $\varepsilon_1$ to blow up spuriously. Using an
\emph{$s$-channel partner} instead (e.g.\ $p_2$, with $\eta(p_1,p_2)=-s/2\sim E^2$) fixes this:
after the correction, $F^1{:}F^2\sim s^{1.0000}$ and $(F^1{:}F^2)(F^3{:}F^4)\sim s^{2.0000}$,
matching naive counting and CGGHJM's own table (2-derivative structures grow as $s$, 4-derivative
as $s^2$). This is the first instance, in either this note or~\cite{fivephotons}, of a recurring
lesson: in any Regge trajectory, using a ``fixed-channel'' pair as a gauge reference produces
spurious growth unrelated to the physical $S$-matrix, and one must classify growing vs.\ fixed
channels \emph{before} constructing explicit polarizations.
\end{remark}

\section{A naive generalization to \texorpdfstring{$n=5$}{n=5} fails structurally}
\label{sec:naive}

The obvious generalization --- boost two particles ($p_1,p_2$) along opposite null reference
directions while holding three ``spectators'' ($p_3,p_4,p_5$) fixed --- fails for structural
reasons, not a computational obstacle. There are two related problems, the second a sharper
version of the same underlying obstruction:

\begin{enumerate}
\item By momentum conservation, if $p_1+p_2$ grows, at least one spectator must absorb the recoil
and grow as well; one cannot hold $p_3+p_4+p_5$ fixed while $p_1+p_2$ diverges, since their sum
must vanish exactly.
\item No generic null momentum can have zero overlap with \emph{both} reference null directions
simultaneously: a null vector with vanishing components along two independent null directions
would live entirely in their spacelike transverse complement, where the only null vector is zero.
Hence invariants like $s_{13}$, which one would naively want to hold fixed (as an analogue of
$t$), necessarily grow too if $p_3$ is a generic null momentum.
\end{enumerate}

At $n=4$, the Regge limit works precisely because all four momenta participate in a specific,
correlated boost-plus-forward-limit structure, not because two particles move and two stay put.
Generalizing this correctly to $n=5$ (where ``holding $t$ fixed'' must mean something precise with
five momenta) is a genuine modeling choice, not resolved in the literature for $n\geq5$, and there
is no evident reason to expect a unique correct choice.

\section{Three families of \texorpdfstring{$n=5$}{n=5} Regge trajectories}
\label{sec:families}

We constructed three qualitatively distinct families, each validated independently (massless
momenta, exact momentum conservation, and an explicit classification of which of the ten pairwise
invariants grow vs.\ stay fixed, checked \emph{before} trusting any polarization construction, per
Remark~\ref{rem:gaugebug}).

\subsection{Family 1: hierarchical hard-trio / recoil-pair}
\label{sec:family1}

A ``hard'' cluster $\{1,2,3\}$ with $s_{12}\to\infty$, and a recoil pair $\{4,5\}$ boosted from a
two-body decay of a system with momentum $P_{45}=-(p_1+p_2+p_3)$. The energy $E_3$ and angle
$\theta(E)$ of $p_3$ relative to $p_1$ are tuned (via the same forward-limit mechanism as
Section~\ref{sec:n4}) to hold \emph{both} $t_{13}:=-(p_1+p_3)^2$ and the recoil invariant mass
$M_{45}^2:=-(p_1+p_2+p_3)^2$ fixed as $E\to\infty$. The direction of the $\{4,5\}\to$ two-body
decay in the $P_{45}$ rest frame cannot be generic (oscillating, no clean power law) or exactly
aligned with the boost axis (degenerate); a small fixed tilt angle $\varphi_0$ gives clean growth
in all ten pairs. Resulting channel classification: growing
$(1,2),(1,4),(1,5),(2,3),(3,4),(3,5)$ (all exactly $\sim E^2$); fixed $(1,3),(2,4),(2,5),(4,5)$
(exact constants).

\subsection{Family 2: two-scale hierarchical}
\label{sec:family2}

Identical construction to Family 1, except the recoil invariant mass is no longer held fixed but
grows at the same rate as the hard channel, $M_{45}^2=m_0^2/x^2$ (writing $x=1/E$). The algebraic
derivation of $\cos\theta,E_3$ that fixes $t_{13}$ and $M_{45}^2$ simultaneously is valid for
\emph{any} value of $M_{45}^2$, including one depending on $x$, so no new derivation was needed.
Resulting channel classification is qualitatively different: only $(1,3)$ (i.e.\ $t_{13}$) stays
fixed; all other nine pairs, including three that were fixed in Family 1, now grow as $\sim E^2$.

\begin{remark}[Relabelings are not new families]
\label{rem:relabel}
Every operator considered here is fully $S_5$-symmetrized (summed over all 120 permutations, as
in~\cite{fivephotons}). Consequently, evaluating such an operator on a trajectory built by
relabeling which trio is ``hard'' and which pair is ``recoil'' gives numerically \emph{identical}
results to the original labeling --- it is not an independent check. A genuine second family
requires a qualitatively different growing/fixed pattern, not merely a different assignment of
particle labels to the same pattern; we verified this obstruction by direct analysis before
building Family 2, and it is the reason Family 2 changes the scaling of $M_{45}^2$ rather than
which particles are ``hard''.
\end{remark}

\subsection{Family 3: fast recoil growth}
\label{sec:family3}

Same construction again, with $M_{45}^2\sim m_0^4/x^4$ (growing \emph{faster} than the hard channel
$s_{12}\sim1/x^2$) --- two powers of $x^{-1}$ beyond Family 2. (A first attempt to instead scale
$t_{13}$ itself at an intermediate rate broke the parity of the kinematics, giving half-integer
powers of $x$ incompatible with our exact-arithmetic machinery; this is consistent with the parity
theorem of~\cite{fivephotons} and was diagnosed before any computation was attempted.) Channel
classification: $(1,3)$ still fixed; $(1,2),(1,4),(1,5)$ grow as before ($\sim E^2$); but
$(2,3),(2,4),(2,5),(3,4),(3,5),(4,5)$ now grow as $\sim E^4$ --- \emph{faster} than the nominally
``hard'' channel itself.

\section{Results}
\label{sec:results}

All computations use exact rational arithmetic (momenta and polarizations constructed in the
quadratic extension $\Q(x)(S)$, $S^2=R(x)$ a rational function --- the only irrationality in the
whole construction, since $S$ always appears linearly until squared --- followed by an exact
Laurent series expansion in $x=1/E$; see~\cite{fivephotons}'s methodology and
Section~\ref{sec:verification} below for the fast backend used). Growth rates are read off as the
lowest surviving power of $x$ in this expansion, i.e.\ $x^{-k}$ corresponds to growth $E^k$.

\subsection{Degree 9}

Within Family 1, the unique degree-9 invariant (within the topologies known before the sandwich
atom was added) grows as $E^8=s^4$, robust across four independent parameter configurations
(exponent $4.00\pm0.01$ throughout). Once the sandwich atom $M^{ij}_{ab}$ is included, the true
dimension of the degree-9 space is 3 (not 1), and an \emph{exact} rank/kernel computation on the
Laurent coefficients (not a numerical extrapolation) proves that $s^3$ is the genuine floor of the
\emph{entire} 3-dimensional space in this trajectory: no combination can do better, confirmed
identically in two independent parameter configurations.

\subsection{Degree 11, Family 1 alone}

Within the (pre-sandwich) 3-dimensional space, an analogous exact argument shows $s^4$ is the
floor: any combination with a component along the two leading directions leaves an $E^9$ residue
that the third cannot cancel. Once the sandwich atom is included systematically, the dimension of
degree 11 jumps to (at least) 15 (documented fully in~\cite{fivephotons}, whose exhaustive
combinatorial sweep later confirmed this is exact within the four known atom types). An exact
rank/kernel computation on this enlarged space reveals a \textbf{9-dimensional subspace} whose
growth is $\leq E^4=s^2$ in Family 1 --- i.e.\ compatible with the conjectured CRG bound ---
identically reproduced in a second independent parameter configuration.

\subsection{Degree 11, two families combined}

The 9-dimensional Family-1 subspace does \emph{not} survive Family 2: all nine directions fall to
$x^{-11}$ ($s^{5.5}$), violating the bound by a wide margin. However, computing the combined kernel
of both families' constraints directly (stacking the Laurent-coefficient rows of both trajectories
and taking the exact nullspace) reveals a genuinely new, smaller \textbf{2-dimensional subspace}
$U_1,U_2$ whose growth is $\leq s^2$ in \emph{both} Family 1 and Family 2 simultaneously ---
verified by direct reconstruction (not just rank counting) and reproduced identically across two
independent parameter configurations in each family. This is the first candidate in this project
to survive more than one qualitatively distinct Regge trajectory.

\subsection{Degree 11, three families combined: a null result}

Family 3 kills $U_1,U_2$ outright: both fall to $x^{-22}$, the same generic order as any
individual, unprotected basis element, violating the bound enormously. A single combination of
$U_1,U_2$ can cancel that leading term, but would then need to independently satisfy roughly 18
more order-by-order constraints using zero remaining degrees of freedom --- there is no structural
reason to expect this, and we did not pursue the (rapidly growing) computational cost of checking
it exhaustively.

Subsequently, a genuinely new atom (the ``triple field-strength sandwich'' $T$) was discovered,
raising the degree-11 dimension from 15 to exactly 16~\cite{fivephotons}. We repeated the
combined-family analysis on this enlarged space. Restricting to Families 1+2, the kernel grows
from dimension 2 to \textbf{dimension 3}: the new direction $U_3$ has a nonzero coefficient of the
$T$-based basis element, confirming it is genuinely new, not a repackaging of $U_1,U_2$. But $U_3$
also falls at $x^{-22}$ in Family 3, exactly like $U_1,U_2$.

To settle this definitively rather than checking individual candidates one at a time, we computed
the combined kernel of \emph{all three} families directly on the complete 16-dimensional space (31
constraint rows: 6 from Family 1, 7 from Family 2, 18 from Family 3, exact rational arithmetic
throughout). The result: \textbf{rank exactly 16, kernel exactly $\{0\}$}. No combination of any of
the sixteen known degree-11 building blocks satisfies $\leq s^2$ in all three trajectory families
simultaneously. This was checked at two independent truncation orders for Family 3 ($N=30,45$),
with identical coefficients across the entire range of orders used --- the numerical stability
required before trusting a high-order Laurent expansion.

\begin{table}[h]
\centering
\small
\begin{tabular}{@{}p{3.2cm} p{2.6cm} p{2.6cm} p{2.6cm}@{}}
\hline
& Family 1 & Families 1+2 & Families 1+2+3 \\
\hline
Degree 9 (dim.\ 3, exact) & floor $s^3$ & --- & --- \\[2pt]
Degree 11, 15 known atoms & subspace dim.\ 9, $\leq s^2$ & subspace dim.\ 2, $\leq s^2$ & \\[2pt]
Degree 11, $+T$ atom (dim.\ 16) & --- & subspace dim.\ 3, $\leq s^2$ & \textbf{kernel $=\{0\}$} \\
\hline
\end{tabular}
\caption{Summary of the degree-9/11 Regge-growth search. Each column to the right imposes the
bound in one additional family simultaneously; the rightmost result is exact and definitive within
the 16-dimensional space explored.}
\end{table}

\section{Verification methodology}
\label{sec:verification}

Beyond the gauge-reference bug of Remark~\ref{rem:gaugebug} and the relabeling non-issue of
Remark~\ref{rem:relabel}, two further real errors were caught and corrected:

\begin{itemize}
\item A first implementation of the fast Laurent-series backend (built after benchmarking showed
sympy's expression-tree overhead, not the underlying arithmetic, was the bottleneck: multiplying
and truncating a series 2000 times took 106s in sympy vs.\ 1.6s with plain \texttt{Fraction}
arithmetic) truncated the high-order end of each series after every multiplication, which is
\emph{incorrect} for Laurent series with negative orders: a discarded high-order term can re-enter
the window of interest after multiplying by a sufficiently negative-order factor later. Caught by a
100-case randomized test against sympy; fixed by never truncating until the final read-out.
\item Re-deriving Family 2 with the fast backend produced mismatches against the (already-trusted)
sympy reference. Diagnosis: the chosen parameters left an uncancelled $\sqrt3$ in the
\emph{individual} (pre-symmetrization) polarizations, which only cancels after summing over the
full $S_5$ orbit; the fast backend's coefficient-conversion routine used Python's \texttt{int()} on
a sympy expression, which silently truncates an irrational like $\sqrt3$ to $1$ instead of raising
an error. Fixed two ways: the conversion routine now requires an exact integer numerator and
denominator and raises an explicit error otherwise; and the Family 2 parameters were changed to
values keeping $\sqrt{R(x)}$ exactly rational throughout, removing the irrationality at its source.
\end{itemize}

As with the companion note, this track record is informative but does not substitute for
independent external verification.

\section{Discussion: what this means}
\label{sec:discussion}

Four structurally different attempts to find a degree-11 candidate compatible with CRG at $n=5$
--- three qualitatively distinct trajectory families, and an enlarged 16-dimensional search space
including a genuinely new atom --- all fail in the same qualitative way: candidates that survive
some subset of the tests are killed by the next, and the combined constraint across everything
tried is $\{0\}$. We think this pattern is more consistent with \emph{the definition of the limit
being wrong} than with \emph{the search being incomplete}, though we cannot rule out the latter
(see the caveats below).

This connects back to the two-problem split made at the start of this project (see the
project notes): deriving CRG from first principles (via the chaos bound, as
in~\cite{chaosbound2021}) was judged out of reach for this project, so we took CRG as a given
conjecture and attacked only its consequences (``Problem B''). What this note shows is that, at
$n\geq5$, Problem B does not cleanly separate from Problem A the way it does at $n=4$: one cannot
even state the bound to be checked without first answering a scaled-down version of the
holographic question --- which kinematic regime is the one in which a causality/chaos-bound
argument would actually constrain growth to $s^2$ at $n\geq5$. That is a real open problem, and we
do not believe it can be resolved by further search within the family of trajectories explored
here.

\begin{remark}[Honest caveats]
This conclusion depends on assumptions we have not independently verified: (i) that the CRG bound,
if it generalizes to $n\geq5$ at all, keeps the same exponent $s^2$ (rather than some
$n$-dependent power); (ii) that the three families explored, despite being qualitatively distinct,
are representative of ``reasonable'' generalizations of the Regge limit, rather than three
correlated failures of a single underlying modeling mistake; (iii) that degree 11 (rather than 13,
15, \ldots) is the right place to look --- the parity theorem of~\cite{fivephotons} rules out even
degrees, but says nothing about whether higher odd degrees behave differently.
\end{remark}

\subsection*{Open problem} What is the correct, physically motivated generalization of the
classical Regge limit to $n\geq5$-particle scattering, and does a version of the Classical Regge
Growth bound hold there? We do not know the answer, and we do not think it can be settled by more
computation of the kind in this note --- it likely requires either a genuine extension of the
holographic chaos-bound derivation of~\cite{chaosbound2021} to $n\geq5$ correlators, or an
independent flat-space causality argument that does not go through holography at all.

\subsection*{Acknowledgments}
The author thanks Claude (Anthropic) for assistance with derivations, computational verification,
and manuscript preparation.

\begin{thebibliography}{9}

\bibitem{cgghjm2019}
S.~D.~Chowdhury, A.~Gadde, T.~Gopalka, I.~Halder, L.~Janagal, S.~Minwalla, \emph{Classifying and
constraining local four photon and four graviton $S$-matrices}, JHEP \textbf{02} (2020) 114;
arXiv:1910.14392.

\bibitem{chaosbound2021}
D.~Chandorkar, S.~D.~Chowdhury, S.~Kundu, S.~Minwalla, \emph{Bounds on Regge growth of flat space
scattering from bounds on chaos}, arXiv:2102.03122.

\bibitem{delduca1993}
V.~Del~Duca, \emph{Parke-Taylor amplitudes in the multi-Regge kinematics}, Phys.\ Rev.\ D
\textbf{48} (1993) 5133; arXiv:hep-ph/9304259.

\bibitem{fivephotons}
M.~Baca, \emph{The multilinear gauge-invariant operator basis for five massless photons:
representation theory, invariant ring, and a complete low-degree classification}, preprint (2026);
\url{https://doi.org/10.5281/zenodo.22819347}.

\end{thebibliography}

\end{document}
